% Sections 2.1-2.6, from lectures/L02/appendix/a_larger_networks.md (A.1-A.6).

\section{Karnaugh maps: from truth table to minimised gate network}
\label{c2:sec:kmap}

Reading a sum-of-products equation straight out of a truth table, the way \secref{c1:sec:boolean}
did, always works but rarely gives the smallest network: every \code{1} row becomes an AND term of
its own, even when several rows share most of their inputs and could be merged.

A \textbf{Karnaugh map} (K-map) makes the shared patterns visible by hand, for functions of 2 to 4
variables, without any formal manipulation of Boolean algebra. A worked example follows; more
practice is in \secref{c2:sec:exercises}, and classical minimisation theory (don't-cares,
Quine\--McCluskey) is outside the scope of this book.

The idea is to lay the truth table out as a grid instead of as a list, and to order the headings so
that every cell differs from each of its neighbours across and down in exactly one input value,
including around the left/right and top/bottom edges. That order is \textbf{Gray code}: for two
bits \code{00, 01, 11, 10}, where only one bit changes at each step, including from the last entry
back to the first.

Then:
\begin{enumerate}
  \item Write a \code{1} in every cell where the output is \code{1}, and leave the rest empty.
  \item Draw rectangular groups around the \code{1}s:
  \begin{itemize}
    \item Every group must be a rectangle whose width and height are each a power of two
      (1, 2 or 4), so that its size is a power of two as well.
    \begin{itemize}
      \item a group of 4 is therefore either a 1x4 line or a 2x2 block, and both are allowed.
      \item cells that meet only at a corner are not a rectangle, so they cannot be grouped.
    \end{itemize}
    \item Groups may overlap, and may wrap around an edge of the grid.
  \end{itemize}
  \item Make every group as large as possible, and cover every \code{1} at least once. A group of 1
    keeps all its inputs in the term, a group of 2 drops one, a group of 4 drops two, and so on.
  \begin{itemize}
    \item Then use as few groups as will do: drop every group all of whose cells are already
      covered by another. Nothing above forbids a redundant group, and a redundant group is a
      redundant AND term in the answer, correct but not minimal.
  \end{itemize}
  \item For every group keep only the input or inputs that have the same value in every cell of the
    group, drop the rest, and write that as an AND term:
  \begin{itemize}
    \item write a kept input plain if the constant value is \code{1}, and primed if it is \code{0}.
    \item a group covering only cells where $B = 0$ therefore contributes the term $B'$.
  \end{itemize}
  \item $X$ is all the groups' terms ORed together.
\end{enumerate}

\subsection*{Worked example}
Three inputs \code{A}, \code{B}, \code{C}, one output \code{X}:

\begin{narrowtable}{cccc}
  \thead{A} & \thead{B} & \thead{C} & \thead{X} \\
  \midrule
  0 & 0 & 0 & 0 \\
  0 & 0 & 1 & 1 \\
  0 & 1 & 0 & 0 \\
  0 & 1 & 1 & 1 \\
  1 & 0 & 0 & 0 \\
  1 & 0 & 1 & 1 \\
  1 & 1 & 0 & 1 \\
  1 & 1 & 1 & 1 \\
\end{narrowtable}

Lay it out with \code{AB} in Gray code down the rows and \code{C} across the columns, and fill in a
\code{1} everywhere $X = 1$:

\lecfig[0.42]{lectures/L02/appendix/images/karnaugh_filled.png}{Karnaugh map for \code{X}, filled in
from the truth table.}{c2:fig:kmap1}

The whole $C = 1$ column is four \code{1}s, a group of 4. Every cell shares $C = 1$ and nothing
else is constant ($A$ and $B$ both vary), so it reduces to the single-variable term $C$:

\lecfig[0.42]{lectures/L02/appendix/images/karnaugh_group_c.png}{The same map with the four-cell
$C = 1$ group marked.}{c2:fig:kmap2}

One \code{1} is still uncovered, on the row $AB = 11$. Group it with its neighbour to the right,
which the $C$ group already covers. Both cells in the pair share $A = 1$ and $B = 1$, and $C$
varies, so it reduces to $AB$:

\lecfig[0.42]{lectures/L02/appendix/images/karnaugh_groups.png}{The same map with the two-cell
$AB = 11$ group added.}{c2:fig:kmap3}

Every \code{1} is now covered; the cell $AB = 11$, $C = 1$ lies in both groups, and overlap is
fine. Summing up gives:

\[ X = AB + C \]

Reading the truth table directly, \chapref{c1:ch}'s way, would have given five AND terms, one per
\code{1} row. The K-map found an equivalent equation with two terms, which is the whole value of
doing this by hand.

\subsection*{Realising the network}
$X = AB + C$ is two gates: an \code{AND} computing $AB$ feeding an \code{OR} whose other input is
$C$.

\lecfig[0.55]{lectures/L02/appendix/images/net.png}{Gate network for the minimised expression.}
{c2:fig:net}

You build this network by hand in CircuitVerse during the session, exactly as you built
\code{or_gate}'s network in \chapref{c1:ch}, and translate it to VHDL in \secref{c2:sec:tovhdl}.

% ------------------------------------------------------------------------------------------
\section{Translating the Karnaugh map's network to VHDL}
\label{c2:sec:tovhdl}

$X = AB + C$ translates straight into VHDL. Three inputs, one output:

\begin{vhdlcode}
library ieee;
use ieee.std_logic_1164.all;

entity gate_network is
    port(a, b, c: in std_logic;
         x      : out std_logic);
end entity;
\end{vhdlcode}

The architecture uses an internal signal \code{ab} to hold the \code{AND} gate's output, and
mirrors the two-gate network one to one:

\begin{vhdlcode}
architecture behaviour of gate_network is
signal ab: std_logic;
begin
    -- AND gate: AB.
    ab <= a and b;

    -- OR gate: AB + C.
    x <= ab or c;
end architecture;
\end{vhdlcode}

Without the intermediate signal it collapses to \code{x <= (a and b) or c;}. Both are equally
valid and synthesise to the same hardware; use whichever is more readable for the network in front
of you.

The buildable version uses the collapsed form, since at two gates there is nothing to gain from
naming the intermediate signal:

\vhdlfile{lectures/L02/gate_network/gate_network.vhd}

Open it beside the version above and convince yourself that they describe the same circuit; that is
the claim this section makes, and it costs nothing to check.

To run it on hardware, assign the ports to switches and an LED with the DE0-CV's pin planner
(appendix~\ref{app:quartus}). Nothing here depends on specific pin numbers.

% ------------------------------------------------------------------------------------------
\section{\code{std_logic_vector}: more than one wire}
\label{c2:sec:vectors}

Everything so far has been \code{std_logic}, a wire carrying one bit. Bundling several wires under
one name is what \code{std_logic_vector} is for.

\begin{vhdlcode}
signal count: std_logic_vector(3 downto 0);
\end{vhdlcode}

That declares four wires: \code{count(3)} down to \code{count(0)}. The range is \code{3 downto 0}
rather than \code{0 to 3} so that the leftmost bit is the most significant, the way you would write
the number on paper. The book uses \code{downto} everywhere, and so does nearly all the VHDL you
will meet.

A vector is a port type like any other:

\begin{vhdlcode}
entity display is
    port(number: in  std_logic_vector(3 downto 0);
         hex   : out std_logic_vector(6 downto 0));
end entity;
\end{vhdlcode}

\subsection*{Values}
A single bit takes single quotes, a vector double ones:

\begin{vhdlcode}
x     <= '1';                -- std_logic: one bit
count <= "1010";             -- std_logic_vector(3 downto 0): four bits
count <= (others => '0');    -- every bit '0', whatever the width
\end{vhdlcode}

The last is an \emph{aggregate}: \code{(others => '0')} means ``all remaining bits are
\code{'0'}'', and since none was named individually, that is all of them. It keeps working when the
width changes, and it turns up throughout the rest of the book.

\subsection*{Indexing and slicing}
One bit out of a vector is a \code{std_logic}; a contiguous run is a \emph{slice}, which is itself a
vector:

\begin{vhdlcode}
carry     <= count(3);
high_half <= input(7 downto 4);
\end{vhdlcode}

Slicing is how one wide port feeds several narrower ones, which is what
\secref{c2:sec:submodules}'s two-digit display does with its eight switch inputs.

One detail is not what a software developer expects. \code{input(7 downto 4)} has the index range
\code{7 downto 4}, not \code{3 downto 0}. Connected to a port declared
\code{std_logic_vector(3 downto 0)}, the two are matched \textbf{positionally, left to right}, not
by index number: \code{input(7)} drives bit 3, \code{input(6)} drives bit 2, and so on. The lengths
have to agree. The numbering does not.

\subsection*{A vector is a bundle of bits, not a number}
\code{count <= "1010"} does not make \code{count} equal to ten. \code{std_logic_vector} says
nothing about what the bits \emph{mean}, so it offers only the operations that make sense for a
bundle of wires: the logical operators applied bit by bit (\code{a and b}, \code{not a},
\code{a xor b}), comparison (\code{count = "1010"}) and concatenation (\code{"10" & "11"} is
\code{"1011"}).

Arithmetic is not on that list. \code{count + 1} will not compile, because nothing has said whether
the four wires are an unsigned number, a signed one, or four unrelated signals.

The book does not work around that with a numeric vector type. Anything that counts is an ordinary
integer with a range:

\begin{vhdlcode}
signal count: natural range 0 to 9;
\end{vhdlcode}

which you can add to and compare directly, and which synthesis turns into exactly the flip-flops the
range needs. Vectors stay what they are: bundles of wires, for things that really are bundles of
wires, such as eight switches or seven display segments.

The two worlds meet only at a boundary, and the conversions turn up where they are needed:

\begin{vhdlcode}
count_v <= std_logic_vector(to_unsigned(count, 4));   -- integer to vector
count   <= to_integer(unsigned(count_v));             -- vector to integer
\end{vhdlcode}

Both come from \code{ieee.numeric_std}, which is why some files add a third \code{use} clause.

% ------------------------------------------------------------------------------------------
\section{Multiplexers}
\label{c2:sec:mux}

A \textbf{multiplexer} (``mux'') connects exactly one of several data inputs to a single output,
chosen by a separate set of \textbf{select inputs}. An \code{N}-to-1 mux has \code{N} data inputs
and $\lceil \log_2 N \rceil$ select bits.

The simplest case, a 2-to-1 mux with the inputs \code{A}, \code{B}, the selector \code{S} and the
output \code{X}:

\begin{narrowtable}{cc}
  \thead{S} & \thead{X} \\
  \midrule
  0 & B \\
  1 & A \\
\end{narrowtable}

As a Boolean equation, $X = S'B + SA$: for every combination of inputs exactly one AND term is
active, and that term's data input passes through.

Larger muxes are the same pattern with more select bits. An 8-to-1 mux with the data inputs
\code{A}--\code{H} and the selector \code{S[2:0]}:

\begin{narrowtable}{rc}
  \thead{S[2:0]} & \thead{X} \\
  \midrule
  000 & A \\
  001 & B \\
  010 & C \\
  011 & D \\
  100 & E \\
  101 & F \\
  110 & G \\
  111 & H \\
\end{narrowtable}

Read as a sum of products with one term per row, the row \code{000} contributes $A\,S_2'S_1'S_0'$
and the rest follow the same shape, each gating its data input with the select combination that
chooses it:

\[ X = A\,S_2'S_1'S_0' + B\,S_2'S_1'S_0 + \dots + H\,S_2S_1S_0 \]

Worth writing out once by hand, not worth memorising. The pattern is the point, and it is the same
at any size.

As gates: one four-input \code{AND} per data input (the data bit plus all three select bits), with
the select bits inverted as required, and all eight \code{AND} outputs summed by a tree of
\code{OR} gates.

\lecfig[0.72]{lectures/L02/appendix/images/mux_8to1.png}{8-to-1 multiplexer as eight AND gates
feeding an OR tree.}{c2:fig:mux8}

This is a real building block. The ATmega328P (the chip on an Arduino Uno and Nano) uses one
internally so that its analogue-capable pins, plus an internal bandgap reference and a temperature
sensor, can share a single analogue-to-digital converter, selected by the four \code{MUX} bits in
its \code{ADMUX} register.

% ------------------------------------------------------------------------------------------
\section{Multiplexers in VHDL: \code{process} and \code{case}}
\label{c2:sec:process}

Every statement so far has been \textbf{concurrent}: a standing description of a wire, all of them
live at once. A multiplexer is the first circuit that is awkward to write that way, and it needs
two new constructs:
\begin{itemize}
  \item a \textbf{\code{process}}, a block whose contents run sequentially, top to bottom, like
    ordinary code. It reruns every time any signal in its \textbf{sensitivity list} changes, and
    from the architecture's point of view it is still just one more concurrent statement. Every
    process in the book has this shape; \chapref{c3:ch} reuses it for clocked logic, with the clock
    in the sensitivity list instead of the inputs being decoded.
  \item a \textbf{\code{case} statement}, VHDL's \code{switch}. Like \code{if} and \code{for} it is
    a \emph{sequential} statement, allowed only inside a process, and every value the selector can
    take must be covered, which \code{when others} guarantees.
\end{itemize}

``Choose one of several inputs based on a selector'' \emph{is} a \code{case} statement, almost word
for word. The 2-to-1 mux from \secref{c2:sec:mux}:

\begin{vhdlcode}
MUX_PROCESS: process(a, b, sel) is
begin
    case (sel) is
        when '0'    => x <= b;
        when '1'    => x <= a;
        when others => x <= '0';
    end case;
end process;
\end{vhdlcode}

Two things to note:
\begin{itemize}
  \item Single-bit values take single quotes, \code{'0'}/\code{'1'}, unlike the double quotes used
    for vectors such as \code{"000"}.
  \item \code{when others} is mandatory rather than defensive. VHDL requires a \code{case} to cover
    every value of the selector's type, and \code{std_logic} has nine (\secref{c1:sec:vhdl}), so
    \code{'0'} and \code{'1'} alone never make a \code{case} exhaustive. It is a language rule, not
    a claim about the wire: a synthesised network carries a \code{0} or a \code{1} and nothing else,
    and the other seven values are the simulator's, where a signal left sitting at \code{'U'} or
    \code{'X'} is a bug rather than a level.
\end{itemize}

\subsection*{Worked example: the 8-to-1 mux}
The same shape with a three-bit selector and eight branches, one per row of \secref{c2:sec:mux}'s
truth table. The data inputs \code{A}--\code{H} become the vector \code{inputs(7 downto 0)}, with
\code{A} as \code{inputs(7)} down to \code{H} as \code{inputs(0)}, so the selector \code{"000"}
connects \code{inputs(7)} through and \code{"111"} connects \code{inputs(0)}.

\vhdlfile{lectures/L02/mux_8to1/mux_8to1.vhd}

Compare it with \secref{c2:sec:mux}'s eight-term equation for the same circuit. The equation
describes the gates; the \code{case} statement describes the \emph{intent} and lets the synthesis
tool produce the gates. Both are correct and both synthesise to the same hardware. That gap,
between describing structure and describing behaviour, is most of what a hardware description
language gives you.

% ------------------------------------------------------------------------------------------
\section{Building a design out of submodules}
\label{c2:sec:submodules}

Everything so far has been one entity with one architecture. That stops being enough the moment a
design needs the same block twice.

The DE0-CV has six 7-segment displays, and you want to drive two of them, each showing one
hexadecimal digit. The logic that turns a four-bit number into seven segment signals is the same
for both. You could write it twice inside one architecture and nothing would stop you, but the
second copy is a liability: every correction has to be made in both, and the day they drift apart
is the day you lose an evening working out which of them is wrong.

Write the block once, as an entity of its own, and \emph{instantiate} it twice.

\subsection*{The submodule}
An ordinary module. Nothing about it says ``I am a submodule'':

\begin{vhdlcode}
entity display is
    port(number: in  std_logic_vector(3 downto 0);
         hex   : out std_logic_vector(6 downto 0));
end entity;
\end{vhdlcode}

Whether it ends up as a whole design or as part of a larger one is decided by whoever instantiates
it, not by the module.

\subsection*{Instantiating it}
With \code{entity work.<name>}, inside another architecture:

\begin{vhdlcode}
entity hex_display is
    port(input     : in  std_logic_vector(7 downto 0);
         hex1, hex0: out std_logic_vector(6 downto 0));
end entity;

architecture structure of hex_display is
begin
    display1: entity work.display
        port map(input(7 downto 4), hex1);

    display0: entity work.display
        port map(input(3 downto 0), hex0);
end architecture;
\end{vhdlcode}

\begin{itemize}
  \item \code{display1} and \code{display0} are \textbf{instance labels}, naming the two copies.
    They have to be unique within the architecture, and they are not the module's name: both
    instances are \code{display}.
  \item \code{entity work.display} selects the module. \code{work} is the library your analysed
    files land in by default, in both GHDL and Quartus.
  \item \code{port map(...)} connects the instance's ports to signals in the enclosing
    architecture, \textbf{in the order the entity declares them}. \code{display} declares
    \code{number} and then \code{hex}, so the first item drives \code{number} and the second is
    driven by \code{hex}.
  \item \code{input(7 downto 4)} is a \textbf{slice}, which the instance sees as a
    \code{std_logic_vector(3 downto 0)}.
\end{itemize}

Note what \code{hex_display}'s architecture does \emph{not} contain: any logic at all. It is wiring
and nothing else. Every gate in the finished design comes from the two \code{display} instances.

\subsection*{Positional and named association}
The \code{port map} above is \textbf{positional} and matches ports by order. VHDL also allows
\textbf{named} association, which you will meet in other people's code:

\begin{vhdlcode}
display1: entity work.display
    port map(number => input(7 downto 4), hex => hex1);
\end{vhdlcode}

The two are equivalent. The book uses positional everywhere, which means the order of an entity's
port declarations is something you have to get right. That is no formality: it is exactly what the
testbenches handed out rely on, which is why every exercise's self-check box spells out its ports
\emph{in order}.

\subsection*{What this costs in hardware}
Nothing surprising. Instantiating a module twice puts two copies of its gates on the FPGA, exactly
as writing the logic twice would. The saving is in the source, not in the silicon: one description,
one place to correct, and a design you can read.
